\subsection{Criterio de descomposición y código de cuentas}
Synaptix adopta una descomposición por entregables. El código raíz \texttt{1} identifica el contrato; los códigos \texttt{1.n} agrupan entregables principales; y los niveles inferiores identifican componentes y paquetes. La etapa contractual será un atributo del trabajo, de modo que una capacidad desarrollada progresivamente conserve su relación con el mismo entregable. La EDT expresa pertenencia y cobertura; las fechas y dependencias se establecerán en la red de actividades.

La descomposición se detendrá cuando el paquete tenga un resultado medible, una condición de aceptación, un responsable único y una estimación defendible. Las instancias actualmente estimadas tienen entre 8 y 80 horas-persona; este rango describe la estimación y no constituye una exigencia de las Bases ni un límite para definir nuevos paquetes. Se mantiene un período quincenal de control. La dedicación parcial y las esperas externas se declararán por separado; el esfuerzo del paquete no equivale automáticamente a su duración. Los servicios recurrentes se presentan mediante familias con resultado, responsable, frecuencia y aceptación comunes; cada período conserva su instancia estimada y asignada.

Cada rama incluirá todo el trabajo de su componente superior y excluirá el trabajo imputado a otra rama. La verificación de cobertura cruzará requisitos, entregables y paquetes; la suma de porcentajes de un gráfico no sustituirá esa comprobación. Las pruebas integradas, los servicios compartidos y las innovaciones tendrán una imputación única; el uso de un componente por varios dominios se registrará como dependencia.

\subsection{Inventario de entregables principales}
El inventario organiza los productos y servicios establecidos en los Subdocumentos 3, 4, 5 y 13. Las Tablas~\ref{tab:sd7-inventario}, \ref{tab:sd7-inventario-datos} y \ref{tab:sd7-inventario-servicio} identifican las ramas y su frontera de trabajo. Sus códigos agrupan entregables y no representan todavía paquetes terminales estimados.

\begin{table}[htbp]
\centering
\begingroup
\setlength{\tabcolsep}{4pt}
\begin{tabularx}{\linewidth}{L{12mm}L{36mm}Y}
\hline
\EncabezadoTabla{Código} & \EncabezadoTabla{Entregable} & \EncabezadoTabla{Cobertura y frontera}\\\hline
1.1 & Expediente de dirección y control & Línea base, trazabilidad, actas, cambios, riesgos, informes y control físico y financiero. Incluye el plan de reversibilidad de los primeros 90 días y sus actualizaciones.\\\hline
1.2 & Plataforma híbrida habilitada & Cinco ambientes DEV, QA, PREPROD, PROD y DR; nube, recinto y terminales; configuración, conectividad, continuidad, despliegue y observabilidad técnica. SD4 y T-11 delimitan los activos.\\\hline
1.3 & Controles de auditoría y seguridad habilitados & Identidades, permisos, consentimiento, cifrado, gestión de claves, trazabilidad y controles transversales. Los ensayos integrales de estos controles pertenecen a 1.10.\\\hline
1.4 & Capacidades del Núcleo Operacional & Operación marítima, Escuela náutica e Infraestructura: salidas y regresos, amarras, participación y retiro, inventario físico e inspección de boyas sin conexión.\\\hline
\end{tabularx}
\caption{Entregables de gobierno, plataforma y núcleo operacional}
\label{tab:sd7-inventario}
\FuenteTabla{Síntesis de Synaptix: SD3, secciones 3.1--3.4; SD4, secciones 4.1--4.3; SD5, secciones 5.1--5.4; SD13, innovaciones 1--5.}
\endgroup
\end{table}

\begin{table}[htbp]
\centering
\begingroup
\setlength{\tabcolsep}{4pt}
\begin{tabularx}{\linewidth}{L{12mm}L{36mm}Y}
\hline
\EncabezadoTabla{Código} & \EncabezadoTabla{Entregable} & \EncabezadoTabla{Cobertura y frontera}\\\hline
1.5 & Capacidades del Núcleo de Servicios & Varadero y contratistas, Combustible y Ambiente y consumos: faenas, habilitación, despacho, medición individual, residuos e indicadores. La medición inicial corresponde a E1; las capacidades completas, a E2.\\\hline
1.6 & Capacidades del Núcleo Administrativo & Clientes y contratos y Finanzas y cobranza: maestros, asignaciones, espera, visitas, hechos facturables, saldos y expedientes. Conserva la autoridad tributaria del sistema contable externo.\\\hline
1.7 & Canales y servicios comunes integrados & Aplicación Android, portal y consola; interfaces externas, notificaciones, telemetría y sincronización. La construcción compartida se imputa aquí; las reglas propias de cada dominio, a 1.4--1.6.\\\hline
1.8 & Datos migrados y archivo histórico consultable & Diagnóstico, herramientas, saneamiento, cargas, conciliación y consulta independiente del sistema de 2014. Comprende la carga inicial de E1 y la cobertura histórica completa de E2 conforme a SD5.\\\hline
\end{tabularx}
\caption{Entregables de servicios, administración, integración y datos}
\label{tab:sd7-inventario-datos}
\FuenteTabla{Síntesis de Synaptix: SD3, secciones 3.1--3.4; SD4, secciones 4.1--4.3; SD5, secciones 5.1--5.4; SD13, innovaciones 1--5.}
\endgroup
\end{table}

\begin{table}[htbp]
\centering
\begingroup
\setlength{\tabcolsep}{4pt}
\begin{tabularx}{\linewidth}{L{12mm}L{36mm}Y}
\hline
\EncabezadoTabla{Código} & \EncabezadoTabla{Entregable} & \EncabezadoTabla{Cobertura y frontera}\\\hline
1.9 & Cinco innovaciones incorporadas y evaluadas & Solicitudes de mantenimiento, mantenimiento preventivo en invernada, lectura asistida, servicio gestionado por resultados y Pasaporte Ambiental Náutico. Incluye su trabajo incremental sin repetir capacidades base.\\\hline
1.10 & Expediente de verificación y aceptación & Protocolos y pruebas integradas, seguridad, desempeño, estrés, aislamiento, recuperación y migración; incidencias y evidencia para T-17. Incluye los dos ensayos completos de migración de SD5.\\\hline
1.11 & Usuarios preparados y etapas implantadas & Materiales por perfil, formación, olas por zona o proceso, convivencia, conciliación diaria, marcha blanca, reversión y estabilización con presencia en terreno.\\\hline
1.12 & Servicio gestionado y transferencia de salida & Operación de ambas etapas en meses 21--56: soporte, NOC/SOC, mantenimiento, respaldos, seguimiento de niveles de servicio y transferencia final de activos, datos y conocimiento.\\\hline
\end{tabularx}
\caption{Entregables de innovación, aceptación, implantación y operación}
\label{tab:sd7-inventario-servicio}
\FuenteTabla{Síntesis de Synaptix: SD3, secciones 3.1--3.4; SD4, secciones 4.1--4.3; SD5, secciones 5.1--5.4; SD13, innovaciones 1--5.}
\endgroup
\end{table}

El inventario separa la construcción de capacidades de su comprobación integral. Cada paquete de construcción incluirá sus pruebas unitarias y documentación; 1.10 reunirá los ensayos entre componentes y la evidencia de recepción. Las pruebas internas específicas de una innovación permanecerán en 1.9 y sus resultados alimentarán el expediente general, sin volver a estimar la misma ejecución.

Las compras y obras físicas atribuidas al cliente por el Subdocumento 2 se registrarán como habilitaciones externas con responsable, fecha requerida y condiciones de recepción. Synaptix mantendrá en sus paquetes las especificaciones, la coordinación y la verificación de integración que le corresponden. Estas dependencias no se tratarán como instalaciones ya ejecutadas ni como ampliaciones automáticas del plazo \parencite[cap.~11; cap.~17, secc.~17.5]{caso07}.

Los códigos \texttt{IN1} a \texttt{IN5} del Subdocumento 13 conservarán su identidad como referencias de origen. El trabajo de diseño, construcción, pruebas y pilotos de las innovaciones se imputará a 1.9; el servicio recurrente, a 1.12. En particular, los informes \texttt{IN1.7.m.q} y la operación \texttt{IN2.7} y \texttt{IN3.7} conservarán su referencia a SD13 dentro de la rama de servicio. Cada paquete de ejecución tendrá una sola imputación de esfuerzo; si un componente de SD13 requiere subdivisión, la matriz de correspondencia identificará todos sus paquetes hijos sin contar también el esfuerzo del padre.

El calendario de la innovación 2 conservará $M_{\mathrm{inv}}$ como el primer junio disponible con IN1 y sus interfaces aceptadas, no anterior al mes 21, y mantendrá la secuencia de invernada y botaduras definida en SD13. Los componentes \texttt{IN4.4} e \texttt{IN5.4} distinguirán sus resultados de puesta en servicio de sus trabajos recurrentes. La correspondencia detallada y las cargas se consolidarán antes de cerrar T-14 y T-15.

\subsection{Expediente de diagnóstico de las fuentes de migración}
El componente \texttt{1.8.1} entregará un expediente que permita decidir cómo extraer, interpretar y conciliar los antecedentes antes de construir las cargas definitivas. Su cobertura comprende los seis conjuntos obligatorios de SD5, sección 5.3.1: socios y armadores; contratos vigentes y documentos; seis años de pagos y saldos; embarcaciones e invernadas; los 14 expedientes de abandono; y tres años de matrícula escolar. El inventario abarcará las fuentes de los doce años de historia para delimitar su tratamiento, sin extender automáticamente a todo ese período la carga operacional \parencite[cap.~15, RT-05.15]{caso07}.

La Figura~\ref{fig:sd7-rama-diagnostico} muestra la pertenencia de los cuatro paquetes al componente 1.8.1. La Tabla~\ref{tab:sd7-paquetes-diagnostico} resume sus resultados y esfuerzos. Las horas son una estimación ascendente inicial de la oferta, basada en las tareas enumeradas en el diccionario; no son mediciones de productividad ni duración comprometida.

\begin{figure}[htbp]
\centering
\begin{tikzpicture}[
 paquete/.style={draw=SynLine,fill=SynSurface,rounded corners=1mm,
 text width=82mm,minimum height=13mm,inner sep=3mm,align=left,
 font=\fontsize{9}{12}\selectfont},
 rama/.style={draw=SynNavy,fill=SynNavy,text=SynWhite,
 text width=32mm,inner sep=3mm,align=left,font=\fontsize{9}{12}\selectfont}]
\node[rama] (raiz) at (-5,0) {\textbf{1.8.1}\\Expediente de diagnóstico de las fuentes de migración};
\node[paquete] (p1) at (2,2.7) {\textbf{1.8.1.1}\\Catálogo de fuentes y accesos};
\node[paquete] (p2) at (2,.9) {\textbf{1.8.1.2}\\Procedimiento de extracción diagnóstica};
\node[paquete] (p3) at (2,-.9) {\textbf{1.8.1.3}\\Informe reproducible de calidad del origen};
\node[paquete] (p4) at (2,-2.7) {\textbf{1.8.1.4}\\Matriz de cobertura y decisiones de migración};
\draw[SynBlue,line width=.8pt] (raiz.east) -- (-2.6,0);
\draw[SynBlue,line width=.8pt] (-2.6,2.7) -- (-2.6,-2.7);
\foreach \paquete/\altura in {p1/2.7,p2/.9,p3/-.9,p4/-2.7}{
 \draw[SynBlue,line width=.8pt] (-2.6,\altura) -- (\paquete.west);
}
\end{tikzpicture}
\caption{Descomposición del expediente de diagnóstico en paquetes}
\label{fig:sd7-rama-diagnostico}
\FuenteFigura{EDT propuesta por Synaptix para el diagnóstico previo a la migración.}
\end{figure}

La jerarquía asigna cada paquete a un único componente superior. Las líneas expresan esa pertenencia; las precedencias de ejecución se definirán separadamente.

\begin{table}[htbp]
\centering
\begin{tabularx}{\linewidth}{L{25mm}YR{13mm}L{31mm}}
\hline
\EncabezadoTabla{Código} & \EncabezadoTabla{Resultado} & \EncabezadoTabla{Horas} & \EncabezadoTabla{Responsable}\\\hline
1.8.1.1 & Catálogo de fuentes y accesos & 24 & Líder de Datos\\\hline
1.8.1.2 & Procedimiento de extracción diagnóstica & 40 & Líder de Datos\\\hline
1.8.1.3 & Informe reproducible de calidad del origen & 32 & Líder de Datos\\\hline
1.8.1.4 & Matriz de cobertura y decisiones de migración & 24 & Líder de Datos\\\hline
\textbf{Total} & \textbf{Expediente de diagnóstico} & \textbf{120} & \\\hline
\end{tabularx}
\caption{Paquetes del expediente de diagnóstico de fuentes}
\label{tab:sd7-paquetes-diagnostico}
\FuenteTabla{Estimación ascendente propuesta por Synaptix; alcance definido en SD5, sección 5.3.}
\end{table}

Las 120 horas comprenden 24 de Análisis Funcional, 72 de Datos, 8 de Seguridad y 16 de Calidad. Incluyen documentación y comprobación de los productos de diagnóstico. No incluyen saneamiento masivo, carga definitiva, consulta histórica implementada, los dos ensayos completos de migración ni la operación del servicio; esos trabajos requieren otros componentes de 1.8, 1.10 y 1.12. La aceptación de este expediente no equivale a la aceptación de la migración.

\ifdefined\EnFormularioTCatorce
El diccionario detallado de los cuatro paquetes se presenta en el apartado siguiente de este formulario.
\else
El diccionario detallado de los cuatro paquetes se presenta en el archivo independiente \path{SYNAPTIX-Formulario-T-14.pdf}.
\fi
Sus fichas conservan estos códigos, esfuerzos, responsables y fronteras; el resumen no sustituye sus condiciones de aceptación.
